% =====================================================================
%  PRÉAMBULE COMMUN — Carnet d'entraînement, Fiche 2 (Nomenclature)
%  (inséré physiquement dans chaque document pour qu'il soit autonome)
% =====================================================================
\documentclass[11pt,a4paper]{article}

% ---------- Encodage & langue ----------
\usepackage[utf8]{inputenc}
\usepackage[T1]{fontenc}
\usepackage{lmodern}
\usepackage[french]{babel}

% ---------- Mathématiques & chimie ----------
\usepackage{amsmath,amssymb,mathtools}
\usepackage{icomma}
\usepackage[version=4]{mhchem}
\usepackage{siunitx}
\sisetup{output-decimal-marker={,},per-mode=symbol}

% ---------- Graphismes & mise en page ----------
\usepackage{xcolor}
\usepackage{colortbl}
\usepackage{tikz}
\usepackage[most]{tcolorbox}
\usepackage{enumitem}
\usepackage{array}
\usepackage{booktabs}
\usepackage{multicol}
\usepackage{fancyhdr}
\usepackage{titlesec}
\usepackage[hidelinks]{hyperref}
\usetikzlibrary{arrows.meta,positioning,calc,shapes.geometric,shapes.misc,
  fit,backgrounds,decorations.pathreplacing}
\usepackage[a4paper,margin=2.2cm,headheight=28pt,headsep=10pt]{geometry}

% =====================================================================
%  PALETTE (charte « Chimie » — mauve/prune du cours)
% =====================================================================
\definecolor{primary}{RGB}{146,95,161}
\definecolor{primarydark}{RGB}{92,55,108}
\definecolor{defcol}{RGB}{0,114,178}
\definecolor{thmcol}{RGB}{0,140,120}
\definecolor{formulacol}{RGB}{198,93,9}
\definecolor{remarkcol}{RGB}{120,94,200}
\definecolor{attncol}{RGB}{200,40,55}
\definecolor{excol}{RGB}{0,135,90}
\definecolor{methcol}{RGB}{90,90,100}
\definecolor{metalcol}{RGB}{198,150,40}
\definecolor{nonmetcol}{RGB}{0,140,150}
\definecolor{oxycol}{RGB}{210,55,45}
\definecolor{hydcol}{RGB}{60,120,210}
\definecolor{catcol}{RGB}{200,70,120}
\definecolor{ancol}{RGB}{40,110,190}
\definecolor{saltcol}{RGB}{120,90,160}

% =====================================================================
%  STYLE COMMUN DES BOÎTES
% =====================================================================
\tcbset{
  boxbase/.style={enhanced,breakable,boxrule=0.9pt,arc=2.5pt,
     left=3mm,right=3mm,top=2.3mm,bottom=2.3mm,
     fonttitle=\bfseries,coltitle=white,
     toptitle=0.6mm,bottomtitle=0.6mm},
}

\newtcolorbox{definition}[1][]{%
  boxbase,colback=defcol!5,colframe=defcol,
  title={\textsc{Définition}\ifx\relax#1\relax\relax\else\textmd{ --- \itshape #1}\fi}}
\newtcolorbox{regle}[1][]{%
  boxbase,colback=thmcol!6,colframe=thmcol,
  title={\textsc{Règle}\ifx\relax#1\relax\relax\else\textmd{ --- \itshape #1}\fi}}
\newtcolorbox{formule}[1][]{%
  boxbase,colback=formulacol!6,colframe=formulacol,
  title={\textsc{Méthode-clé}\ifx\relax#1\relax\relax\else\textmd{ --- \itshape #1}\fi}}
\newtcolorbox{remarque}[1][]{%
  boxbase,colback=remarkcol!6,colframe=remarkcol,
  title={\textsc{Remarque}\ifx\relax#1\relax\relax\else\textmd{ : \itshape #1}\fi}}
\newtcolorbox{attention}[1][]{%
  boxbase,colback=attncol!6,colframe=attncol,
  title={\textsc{Attention}\ifx\relax#1\relax\relax\else\textmd{ : \itshape #1}\fi}}
\newtcolorbox{exemple}[1][]{%
  boxbase,colback=excol!5,colframe=excol,
  title={\textsc{Exemple}\ifx\relax#1\relax\relax\else\textmd{ : \itshape #1}\fi}}
\newtcolorbox{methode}[1][]{%
  boxbase,colback=methcol!7,colframe=methcol,sharp corners,
  title={\textsc{Méthode}\ifx\relax#1\relax\relax\else\textmd{ : \itshape #1}\fi}}
\newtcolorbox{astuce}[1][]{%
  boxbase,colback=primary!6,colframe=primary,
  title={\textsc{Astuce concours}\ifx\relax#1\relax\relax\else\textmd{ : \itshape #1}\fi}}

% ---- Exercice (numéroté) ----
\newtcolorbox[auto counter]{exercice}[1][]{%
  boxbase,colback=primary!4,colframe=primarydark,
  title={\textsc{Exercice}~\thetcbcounter\ifx\relax#1\relax\relax\else\textmd{ --- \itshape #1}\fi}}

% ---- Corrigé (numéroté) ----
\newtcolorbox[auto counter]{corrige}[1][]{%
  boxbase,colback=excol!4,colframe=excol,
  title={\textsc{Corrigé}~\thetcbcounter\ifx\relax#1\relax\relax\else\textmd{ --- \itshape #1}\fi}}

% =====================================================================
%  EN-TÊTES ET PIEDS DE PAGE
% =====================================================================
\pagestyle{fancy}
\fancyhf{}
\fancyhead[L]{\small\textcolor{primarydark}{\textbf{Fiche 2} — Nomenclature}}
\fancyhead[R]{\small\textcolor{primarydark}{\textsc{Carnet d'entraînement}}}
\fancyfoot[C]{\small\thepage}
\renewcommand{\headrulewidth}{0.8pt}
\renewcommand{\headrule}{\hbox to\headwidth{\color{primary}\leaders\hrule height \headrulewidth\hfill}}

\titleformat{\section}{\Large\bfseries\color{primarydark}}{\thesection}{0.6em}{}
\titleformat{\subsection}{\large\bfseries\color{primary}}{\thesubsection}{0.6em}{}
\titleformat{\subsubsection}{\normalsize\bfseries\color{primary!80!black}}{}{0em}{}

% ---- Bandeau de titre réutilisable : \bandeau{Thème}{Sous-titre} ----
\newcommand{\bandeau}[2]{%
  \begin{center}
  \begin{tikzpicture}
  \node[rectangle,rounded corners=7pt,fill=primary,text=white,
        minimum width=16.0cm,minimum height=1.7cm,align=center,inner sep=8pt]
    {{\Large\bfseries #1}\\[3pt]{\normalsize #2}};
  \end{tikzpicture}
  \end{center}
  \vspace{0.3cm}}
\begin{document}
\bandeau{Thème 2 — Difficile}{Du nom à la formule : croisement des valences --- Exercices}

\noindent\textit{On construit la difficulté par sous-questions, en apprenant aussi à \emph{écarter} les
réponses fausses, comme au concours.}
\vspace{0.3cm}

\begin{exercice}[nitrate d'étain(IV) : la bonne formule et les pièges]
\begin{enumerate}[label=\alph*),leftmargin=1.8em,topsep=2pt,itemsep=1.5pt]
  \item Quel est le symbole de l'étain ? Quelle est la charge du cation « étain(IV) » ?
  \item Quelle est la formule et la valence de l'ion nitrate ? En quoi diffère-t-il du nitrite ?
  \item Croise les valences et écris la formule finale (avec parenthèses et vérification de neutralité).
  \item On propose les réponses \ce{Sr(NO3)2}, \ce{Sr(NO3)4}, \ce{Sn(NO3)2}. Explique pourquoi chacune est
        fausse.
\end{enumerate}
\end{exercice}

\begin{exercice}[sulfite d'aluminium : ne pas confondre l'anion]
\begin{enumerate}[label=\alph*),leftmargin=1.8em,topsep=2pt,itemsep=1.5pt]
  \item Donne la formule et la valence de l'ion \emph{sulfite}, puis de l'ion \emph{sulfate}.
  \item Écris la formule du \textbf{sulfite d'aluminium} (croisement, parenthèses, neutralité).
  \item Parmi \ce{Al(SO3)3}, \ce{Al2(SO4)3}, \ce{AlSO4}, laquelle (le cas échéant) pourrait être confondue
        avec ta réponse, et pourquoi est-elle fausse ?
\end{enumerate}
\end{exercice}

\begin{exercice}[la grille des croisements]
Remplis la grille : pour chaque couple (cation, anion), écris la formule brute neutre.
\begin{center}
\renewcommand{\arraystretch}{1.4}
\begin{tabular}{c|ccc}
\toprule
 & \ce{Cl-} & \ce{SO4^2-} & \ce{PO4^3-} \\
\midrule
\ce{Na+}  & \phantom{XXXX} & \phantom{XXXXX} & \phantom{XXXXX} \\
\ce{Ca^2+} & & & \\
\ce{Al^3+} & & & \\
\bottomrule
\end{tabular}
\end{center}
Que remarques-tu sur la diagonale où les valences du cation et de l'anion sont égales ?
\end{exercice}

\begin{exercice}[lecture inverse : de la formule aux ions]
Pour chaque formule, retrouve le \textbf{cation}, l'\textbf{anion} et \emph{vérifie} la neutralité. Déduis-en
le nom du composé.
\begin{enumerate}[label=\alph*),leftmargin=1.8em,topsep=2pt,itemsep=1.5pt]
  \item \ce{Mg(ClO4)2} \qquad b) \ce{(NH4)2CO3} \qquad c) \ce{Fe2(SO4)3} \qquad d) \ce{Ca3(PO4)2}
\end{enumerate}
\end{exercice}

\begin{exercice}[déduire la valence d'un métal à partir de la formule]
On te donne des formules neutres où l'anion est connu ; déduis la \emph{valence} (donc le degré) du métal,
puis nomme-le en précisant « métal(\ldots) ».
\begin{enumerate}[label=\alph*),leftmargin=1.8em,topsep=2pt,itemsep=1.5pt]
  \item \ce{Cu3(PO4)2} : quelle est la valence du cuivre ? \quad (anion \ce{PO4^3-})
  \item \ce{Sn(SO4)2} : quelle est la valence de l'étain ? \quad (anion \ce{SO4^2-})
  \item \ce{Pb(NO3)2} \emph{et} \ce{Pb(NO3)4} : donne la valence du plomb dans chacun.
\end{enumerate}
\end{exercice}
\end{document}
